\documentclass[12pt]{article}
\usepackage[utf8]{inputenc}
\usepackage[T1]{fontenc}
\usepackage{hyperref}
\title{Tejalapan Zapotec Language Resources for K-12 Curriculum Development}
\author{Compiled by Claude Sonnet 5.5 (claude-sonnet-5-5) for Andr\'as Kornai}
\date{}
\begin{document}
\maketitle

\section{Introduction}
Tejalapan Zapotec (ISO 639-3 ztt, Glottocode teja1235) is a Zapotecan language of San Felipe Tejal\'apam, Oaxaca, Mexico. It is severely endangered. Speaker counts range from about 50 to 124, depending on the source. One source describes it as unwritten. Searches found only catalogue-type entries, with no teaching materials.

\section{Primary Educational Resources}
None found. No primers, textbooks or school programmes were located.

\section{Linguistic and Reference Materials}
\begin{itemize}
\item Wikipedia overview: \url{https://en.wikipedia.org/wiki/Tejalapan_Zapotec}
\item Endangered-language profile with speaker estimates (HLT, Budapest): \url{https://hlt.bme.hu/en/dld/language/310}
\item Verbix endangered-languages page: \url{https://docs.verbix.com/EndangeredLanguages/TejalapanZapotec}
\item Glottolog entry: \url{https://glottolog.org/resource/languoid/id/teja1235.bigmap.html}
\end{itemize}

\section{Religious and Community Resources}
No scripture or community organisation was found. A Joshua Project page exists, but its data looks unreliable (it gives a population and alternate name that do not match other sources): \url{https://joshuaproject.net/people_groups/15346/MX}

\section{Digital and Audio Resources}
\begin{itemize}
\item Global Recordings Network listing: \url{https://www.globalrecordings.net/en/language/ztt}
\end{itemize}

\section{Curriculum Development Context}
With very few speakers, mostly older adults, and no known written materials, any curriculum would depend on documentation and revitalisation work with the remaining speakers. The sources found support only a very modest starting point.
\end{document}